\subsection{判别式--谱曲线编译下最近上半平面分支的半代数闭包与复杂度证书}\label{subsec:conclusion-discriminant-widom-semi-algebraic-nash-complexity}
在有限状态有效配分函数的判别式--谱曲线编译框架中，固定截断阶 \(d\) 后可得到特征多项式 \(\Pi_d(\lambda,y;T)\) 及其判别式
\[
D_d(y;T):=\Disc_{\lambda}\!\bigl(\Pi_d(\lambda,y;T)\bigr).
\]
以下结果对一般多项式族
\[
D(y;T)\in\RR[y,T_1,\dots,T_p]
\]
成立，其中 \(T\in\RR^p\) 为实参数，\(y=a+\mathrm{i}b\)（\(a,b\in\RR\)）。记
\[
F(T,a,b):=\Re D(a+\mathrm{i}b;T),\qquad
G(T,a,b):=\Im D(a+\mathrm{i}b;T),
\]
则 \(F,G\in\RR[T,a,b]\)。

\begin{theorem}[最近上半平面分支图形的半代数性闭包]\label{thm:conclusion-widom-nearest-branch-semi-algebraic-closure}
定义上半平面零点集
\[
S:=\{(T,a,b)\in\RR^{p+2}: b>0,\ F(T,a,b)=0,\ G(T,a,b)=0\},
\]
并定义最近上半平面零点图形
\[
W:=\left\{
(T,a,b)\in S:\ \nexists(a',b')\ \text{s.t.}\ (T,a',b')\in S,\ b'<b
\right\}.
\]
则 \(W\subset\RR^{p+2}\) 为半代数集；其参数投影
\[
P:=\pi_T(W)\subset\RR^p
\]
亦为半代数集。特别地，对每个固定 \(d\)，由最近上半平面分支诱导的 Widom 临界对象在参数空间中必为半代数集合。
\end{theorem}
\leanverified{paper\_conclusion\_widom\_nearest\_branch\_semi\_algebraic\_closure}
\begin{proof}
集合 \(S\) 由多项式等式与不等式给出，故为基本半代数集。构造
\[
U:=\{(T,a,b,a',b'):\ (T,a,b)\in S,\ (T,a',b')\in S,\ b'<b\},
\]
则 \(U\) 亦为半代数集。由 Tarski--Seidenberg 定理，投影
\[
V:=\pi_{T,a,b}(U)
\]
半代数。按定义 \(W=S\setminus V\)，而半代数类对布尔运算封闭，故 \(W\) 半代数。再次投影即得 \(P=\pi_T(W)\) 半代数。证毕。
\end{proof}

\begin{theorem}[简单最近分支的 Nash 正则性与失正则证书]\label{thm:conclusion-widom-nearest-branch-nash-regularity}
在定理 \ref{thm:conclusion-widom-nearest-branch-semi-algebraic-closure} 设定下，定义
\[
W^{\mathrm{reg}}
:=
\left\{
(T,a,b)\in W:
\begin{array}{l}
\text{(i) 在纤维 }S_T:=\{(a',b'):(T,a',b')\in S\}\text{ 中，}\\
\qquad\text{最小虚部点唯一；}\\
\text{(ii) }\partial_yD(a+\mathrm{i}b;T)\neq 0
\end{array}
\right\}.
\]
则：
\begin{enumerate}
  \item \(W^{\mathrm{reg}}\) 是 \(W\) 的半代数相对开子集；
  \item 在 \(P^{\mathrm{reg}}:=\pi_T(W^{\mathrm{reg}})\) 的每个连通分支上，\(W^{\mathrm{reg}}\) 是某个映射
  \[
  y_{\ast}(T)=a_{\ast}(T)+\mathrm{i}\,b_{\ast}(T)
  \]
  的图形，且 \(T\mapsto y_{\ast}(T)\) 为 Nash 函数；
  \item 全部失正则参数包含于两类半代数集合之并：
  \[
  \mathcal B_{\mathrm{alg}}
  :=
  \pi_T\!\{(T,a,b):F=G=0,\ \partial_yD(a+\mathrm{i}b;T)=0\},
  \]
  \[
  \mathcal B_{\mathrm{tie}}
  :=
  \left\{
  T:\ \exists (a_1,b),(a_2,b)\in S_T,\ (a_1,b)\neq(a_2,b),\ b=\min_{(a',b')\in S_T}b'
  \right\}.
  \]
\end{enumerate}
\end{theorem}
\begin{proof}
“唯一最近点”与“\(\partial_yD\neq 0\)”均可写成半代数量词公式（后者等价于 \(|\partial_yD|^2>0\)），故 \(W^{\mathrm{reg}}\) 半代数。
由非退化条件 \(\partial_yD\neq 0\)，等价于 \((F,G)\) 对 \((a,b)\) 的 Jacobian 非奇异，故由实解析隐函数定理，局部存在唯一实解析分支 \((a_\ast,b_\ast)\)。
又因该图形来自半代数集合，在连通分支上可拼接为全局半代数解析图形，即 Nash 图形。
失正则只能由代数分歧（导数退化）或最近性并列（最小虚部多点并列）触发，故包含于 \(\mathcal B_{\mathrm{alg}}\cup\mathcal B_{\mathrm{tie}}\)。证毕。
\end{proof}

\begin{corollary}[固定截断阶下 Widom 几何的驯服性]\label{cor:conclusion-widom-tame-geometry-one-parameter}
令 \(p=1\)，并在 \(P^{\mathrm{reg}}\subset\RR\) 上取任一 Nash 最近分支
\[
T\longmapsto b_\ast(T):=\Im y_\ast(T).
\]
则：
\begin{enumerate}
  \item \(b_\ast\) 在其定义域上仅有有限多个不连续点与有限多个局部极值点；
  \item 不存在有限参数区间内的无限振荡与分形边界累积；
  \item 对固定 \(d\)，最近 Lee--Yang/Widom 轨迹与对应临界集的拓扑复杂度有限。
\end{enumerate}
\end{corollary}
\begin{proof}
半代数集合构成 \(o\)-minimal 结构；单变量可定义函数满足分段单调与有限分段光滑。Nash 函数尤其可定义，故极值点与不连续点均为有限个，进而排除无限振荡与分形累积。证毕。
\end{proof}

\begin{proposition}[分支点与不可解析参数的结果式证书]\label{prop:conclusion-widom-resultant-certificate-degree-bound}
\leanverified{paper\_conclusion\_widom\_resultant\_certificate\_degree\_bound}
设 \(p=1\)，且 \(D(y;T)\in\QQ[T][y]\)。记
\[
n:=\deg_y D,\qquad m:=\deg_T D,
\]
并定义
\[
R(T):=\Res_y\!\bigl(D(y;T),\ \partial_yD(y;T)\bigr)\in\QQ[T].
\]
则：
\begin{enumerate}
  \item 若 \(R(T_0)\neq 0\)，则 \(D(\cdot;T_0)\) 无重根，故其 \(n\) 个根可在 \(T_0\) 邻域解析追踪为局部解析分支 \(y_j(T)\)；
  \item 结果式次数满足
  \[
  \deg_T R\le (2n-1)m.
  \]
  因而重根/分歧参数点（计重数）至多 \((2n-1)m\) 个，实轴上亦至多此数量级。
\end{enumerate}
\end{proposition}
\begin{proof}
结果式基本性质给出
\[
R(T_0)=0
\Longleftrightarrow
\gcd_y\!\bigl(D(\cdot;T_0),\partial_yD(\cdot;T_0)\bigr)\neq 1
\Longleftrightarrow
D(\cdot;T_0)\ \text{有重根}.
\]
故 \(R(T_0)\neq 0\) 时根皆简单；由隐函数定理得到局部解析延拓。
次数估计由一般界
\[
\deg_T\Res_y(P,Q)\le (\deg_y Q)\deg_T P+(\deg_y P)\deg_T Q
\]
应用于 \(P=D,\ Q=\partial_yD\) 得
\[
\deg_T R\le (n-1)m+n\,m=(2n-1)m.
\]
证毕。
\end{proof}

\begin{corollary}[最近分支解析片段数的全局有界性]\label{cor:conclusion-widom-event-set-global-finite}
在命题 \ref{prop:conclusion-widom-resultant-certificate-degree-bound} 设定下，定义事件集
\[
\mathcal E:=\{T\in\RR:R(T)=0\}\ \cup\ \mathcal B_{\mathrm{tie}}.
\]
则：
\begin{enumerate}
  \item \(\mathcal E\) 为半代数集，故由有限个点与有限个区间组成；在泛型非退化情形下，\(\mathcal E\) 为有限点集；
  \item \(\RR\setminus\mathcal E\) 分解为有限多个开区间，每个区间上最近分支可取为 Nash 函数，并与某一固定代数根分支 \(y_j(T)\) 一致；
  \item 固定截断阶 \(d\) 时，最近分支的跳变次数、解析片段数与拐点数均有限，且可由 \((n,m)\) 的显式函数控制，其中一部分上界由 \(\deg R\le (2n-1)m\) 直接给出。
\end{enumerate}
\end{corollary}
\begin{proof}
\(\mathcal B_{\mathrm{tie}}\) 由半代数量词条件定义，故半代数；\(\{R=0\}\) 代数，故 \(\mathcal E\) 半代数。
一维半代数集可作有限 cell 分解，得到有限区间结构。
在每个无事件区间内，最近分支唯一且 \(\partial_yD\neq 0\)，由定理 \ref{thm:conclusion-widom-nearest-branch-nash-regularity} 得 Nash 解析图形，并与同一代数分支一致。证毕。
\end{proof}

\leanverified{paper\_conclusion\_widom\_definable\_selection\_global}
\begin{theorem}[最近分支的全局半代数可选取性]\label{thm:conclusion-widom-definable-selection-global}
沿用定理 \ref{thm:conclusion-widom-nearest-branch-semi-algebraic-closure} 的记号，令 \(P=\pi_T(W)\subset\RR^p\)。则存在半代数映射
\[
\sigma:P\to\RR^2,\qquad
\sigma(T)=\bigl(a_\sigma(T),b_\sigma(T)\bigr),
\]
使得
\[
(T,a_\sigma(T),b_\sigma(T))\in W\qquad(\forall\,T\in P).
\]
并且存在 \(P\) 的有限半代数分割 \(\{P_\nu\}\)，满足：
\begin{enumerate}
  \item 在每个 \(P_\nu\) 上，\(\sigma\) 可取连续；
  \item 若 \(P_\nu\) 落在正则区域（\(\partial_yD\neq 0\) 且最近点唯一），则 \(\sigma|_{P_\nu}\) 可取为 Nash。
\end{enumerate}
\end{theorem}
\begin{proof}
由 \(W\) 半代数与可定义选择定理，投影 \(P=\pi_T(W)\) 存在半代数截面 \(\sigma\)。
再由 cell 分解可将 \(P\) 分成有限块，使 \(\sigma\) 在每块上连续。
在正则块上，隐函数定理给出实解析性；结合半代数图形得到 Nash 性。证毕。
\end{proof}

\begin{conjecture}[Puiseux 指数与有限体积最近零点标度的一一对应]\label{conj:conclusion-widom-puiseux-scaling-volume-law}
设 \(p=1\)，并令 \((T_c,y_c)\) 为代数分歧点。若局部归一化满足
\[
y(T)-y_c\sim C\,(T-T_c)^{1/q},\qquad C\neq 0,\ q\ge 2,
\]
则对由有限状态递推（或伴随矩阵）生成的有限体积多项式序列 \(\{Z_{V,d}(y;T)\}_V\)，其最近上半平面零点应满足
\[
\Im y_{V,\ast}(T_c)\asymp V^{-1/q}.
\]
并且对应 Richardson 外推应以阶 \(q\) 自然闭合，存在将误差消去至 \(V^{-(1+1/q)}\) 量级的编译方案。该对应将代数分歧指数与有限体积标度指数直接联结为可审计读取规则。
\end{conjecture}

\begin{remark}[几何复杂度与可计算证书的统一接口]\label{rem:conclusion-widom-semi-algebraic-complexity-interface}
上述定理链将最近上半平面分支的几何与计算复杂度归约为三类对象：半代数闭包、Nash 正则分层与结果式次数证书。由此，固定截断阶下的 Widom/Lee--Yang 最近分支不再依赖无界分岔叙述，而可在有限证书框架中完成全局审计。
\end{remark}
